\chapter{Otimização de \textit{timing} e de área}
\label{cap:otimizacao}
\textcolor{red}{[PENDENTE: parágrafo de abertura do capítulo]}

\section{Metodologia}

O processador foi sintetizado para o FPGA Artix-7 \texttt{xc7a100tcsg324-1}
da placa Nexys A7-100T, com restrição de \textit{clock} de 100~MHz (período
de 10~ns). As medições foram feitas no Vivado 2025.2, num fluxo em lote fixo
composto por síntese (\texttt{synth\_design}), otimização lógica
(\texttt{opt\_design}), posicionamento (\texttt{place\_design}) e roteamento
(\texttt{route\_design}), todos com as diretivas padrão. Todos os valores
são pós-roteamento. A métrica principal é a folga do pior caminho
(\textit{Worst Negative Slack}, WNS). A partir dela, estima-se a frequência
máxima de operação:
\begin{equation}
  F_{\max} = \frac{1}{T - \mathrm{WNS}},
\end{equation}
em que $T$ é o período da restrição. Um WNS negativo indica que o caminho
crítico não cabe no período: com $\mathrm{WNS} = -4$~ns, por exemplo, o
caminho leva 14~ns e $F_{\max} \approx 71{,}4$~MHz.

Três cuidados garantiram que as comparações fossem válidas:
\begin{itemize}
  \item todas as medições de uma comparação usam o mesmo programa na memória
    de instruções, já que o conteúdo dessa memória altera o circuito
    sintetizado (Seção~\ref{sec:programa-teste}); a partir da segunda etapa,
    o programa é o autoteste com as 37 instruções;
  \item cada mudança de RTL só foi medida depois de passar por todos os
    \textit{testbenches}, de modo que nenhuma otimização alterou o
    comportamento do processador sem que isso fosse detectado;
  \item foi alterada uma estrutura por vez, para que o efeito de cada
    mudança pudesse ser isolado.
  \item o fluxo é determinístico: a mesma entrada produz exatamente o mesmo
    resultado. Porém, pequenas alterações no circuito mudam o posicionamento.
    Dois circuitos com a mesma lógica e apenas duas LUTs de diferença
    apresentaram WNS diferentes em 0,47~ns. Por isso, diferenças abaixo de
    cerca de 0,5~ns foram consideradas variação de posicionamento, e não
    efeito do RTL.
\end{itemize}

\section{Diagnóstico inicial}

Na versão inicial, o processador atingia apenas 45,2~MHz
($\mathrm{WNS} = -12{,}1$~ns). O caminho crítico ia do PC, passava pela
memória de instruções, pela decodificação, pela ALU e pela memória de
dados, e terminava na escrita do banco de registradores. Eram 18 níveis de
lógica e 21,8~ns de atraso, dos quais 85\% correspondiam a roteamento, e não
a lógica.

O relatório de utilização apontou a causa: o circuito tinha 32.867
\textit{flip-flops}, um número próximo dos 32.768 bits da memória de dados
de 4~KiB. A memória não havia sido reconhecida como memória pela síntese e
estava implementada com registradores comuns. Além do desperdício de área,
o multiplexador de leitura de 1024 posições ficava espalhado pelo chip. Isso
explica a predominância do roteamento no atraso.

\section{Memória de dados como RAM distribuída}

A memória de dados estava descrita como uma matriz bidimensional de bytes,
com leitura combinacional e escrita por byte:

\begin{lstlisting}
logic [7:0] mem [0:MEM_DEPTH-1][0:3];   // 1024 palavras x 4 bytes
\end{lstlisting}

A síntese do Vivado só infere memória a partir de padrões de descrição
específicos, e essa forma não era reconhecida. A memória foi reescrita como
um vetor unidimensional de palavras de 32 bits, em que a escrita de cada byte
é habilitada separadamente dentro de um único bloco síncrono, sem
\textit{reset} e sem lógica adicional sobre o vetor:

\begin{lstlisting}
logic [31:0] mem [0:MEM_DEPTH-1];

always_ff @(posedge clk)
    if (ram_write_en)
        for (int b = 0; b < 4; b++)
            if (we_byte[b]) mem[word_addr][8*b +: 8] <= wdata[8*b +: 8];
\end{lstlisting}

O arquivo de inicialização continua com um byte por posição. Um bloco
\texttt{initial} lê esse arquivo para um vetor temporário e monta cada palavra
em ordem \textit{little-endian}, de forma que o formato dos arquivos gerados
pelo compilador não mudou.

Os blocos de BRAM do FPGA não foram usados porque têm leitura síncrona, com
um ciclo de latência. Num processador monociclo, o dado lido precisa estar
disponível no mesmo ciclo em que o endereço é calculado. Com a nova
descrição, a memória passou a ser inferida como RAM distribuída (LUTRAM),
implementada nas próprias LUTs do FPGA e com leitura assíncrona. O
comportamento não mudou.

O número de \textit{flip-flops} caiu de 32.867 para 99, e o de LUTs, de
15.204 para 1.685 (redução de 89\%). A frequência estimada subiu de 45,2
para 57,5~MHz (+27\%). O atraso do caminho crítico caiu para 17,1~ns, e o
pior caminho passou a ser o de uma instrução de \textit{load} completa.

\section{Somador dedicado para o endereço de memória}

Na arquitetura original, o endereço dos acessos à memória
($\mathit{rs1} + \mathit{imediato}$) era calculado pela própria ALU.
Assim, o caminho de um \textit{load} passava pela decodificação do sinal que
seleciona o segundo operando, pelo multiplexador desse operando, pelo
somador da ALU e pelo multiplexador que escolhe a operação no resultado da
ALU. Só então o endereço chegava à memória. Foi acrescentado à unidade de
execução um somador exclusivo para o endereço, em paralelo à ALU e ligado
diretamente à memória de dados:

\begin{lstlisting}
assign mem_addr_o = reg_data1_i + imm_i;   // somador dedicado, independe da ALU
\end{lstlisting}

O resultado da ALU continua indo para a etapa de \textit{write-back}, de
forma que o comportamento das instruções não mudou. Com a decodificação do
seletor e os dois multiplexadores fora do caminho crítico, o WNS passou de
$-7{,}6$ para $-4{,}2$~ns, e a frequência estimada, de 56,7 para 70,4~MHz
(+24\%). O custo foi de 117 LUTs.

\section{Entradas e saídas genéricas}

Na primeira versão, as saídas do MMIO podiam ser lidas de volta pelo software. A
medição mostrou que isso piorava o WNS em 0,9~ns e acrescentava 96
\textit{flip-flops}. Como o programa pode ler qualquer bit das saídas, a
síntese não pode remover os bits que não estão ligados a nenhum pino. As
saídas passaram então a ser somente de escrita, e a leitura desses
endereços retorna zero. Nenhum dos programas lia as saídas. Com isso, a
síntese descarta os bits sem pino: em relação à etapa anterior, o
processador ficou com 125 LUTs e 32 \textit{flip-flops} a menos. O WNS
medido, de $-4{,}7$~ns, está dentro da faixa de variação de posicionamento:
uma variante com a mesma lógica obteve $-4{,}3$~ns.

\section{Resultados}

A Tabela~\ref{tab:otimizacao} resume a evolução. As duas primeiras linhas
foram medidas com o programa \texttt{selftest} e as demais com o autoteste
de 37 instruções. A linha de referência mede o mesmo RTL da etapa de
LUTRAM com o novo programa e mostra que os dois programas são equivalentes
para fins de medição.

\begin{table}[htbp]
  \centering
  \small
  \caption{Evolução de \textit{timing} e área (pós-roteamento, restrição de 10~ns).}
  \label{tab:otimizacao}
  \begin{tabular}{@{}lrrrrrrr@{}}
    \toprule
    Etapa & WNS (ns) & $F_{\max}$ (MHz) & Atraso (ns) & Níveis & Rot. & LUTs & FFs \\
    \midrule
    Inicial                     & $-12{,}115$ & 45,2 & 21,8 & 18 & 85\% & 15.204 & 32.867 \\
    Memória em LUTRAM           & $-7{,}391$  & 57,5 & 17,1 & 15 & 81\% & 1.685  & 99 \\
    Referência (autoteste)      & $-7{,}625$  & 56,7 & 17,3 & 18 & 81\% & 1.684  & 99 \\
    Somador dedicado            & $-4{,}207$  & 70,4 & 13,9 & 20 & 70\% & 1.801  & 99 \\
    E/S com leitura de volta\textsuperscript{a} & $-5{,}083$ & 66,3 & 14,8 & 21 & 73\% & 1.798 & 195 \\
    E/S genérica (final)        & $-4{,}724$  & 67,9 & 13,8 & 20 & 72\% & 1.676  & 67 \\
    \bottomrule
    \multicolumn{8}{@{}l}{\footnotesize \textsuperscript{a} Versão descartada.
      Rot.: parcela do atraso devida a roteamento. LUTRAM: 556 LUTs em todas as
      etapas após a segunda.}
  \end{tabular}
\end{table}

Da versão inicial à final, a frequência máxima estimada subiu de 45,2 para
cerca de 68--70~MHz, um ganho de aproximadamente 50\%. O atraso do caminho
crítico caiu de 21,8 para 13,8~ns. A área caiu para 11\% das LUTs e 0,2\%
dos \textit{flip-flops} iniciais. Esses valores são do fluxo fixo de
medição. No fluxo padrão de implementação do projeto no Vivado, que inclui a
etapa de otimização física (\texttt{phys\_opt\_design}), a versão com
somador dedicado atingiu WNS de $-3{,}6$~ns, cerca de 73~MHz. O ganho veio
principalmente da replicação dos registradores do PC, cuja saída alimenta
cerca de 190 entradas.

\section{Caminho crítico atual e limitações}

O caminho crítico atual é o de uma instrução de \textit{load}, com 20 níveis
de lógica e cerca de 70\% do atraso em roteamento. Na versão com somador
dedicado, ele se dividia aproximadamente assim:
\begin{itemize}
  \item busca da instrução na memória de instruções (3,3~ns);
  \item leitura do banco de registradores (2,3~ns);
  \item soma do endereço, com propagação do \textit{carry} até o bit 31
    (2,6~ns);
  \item decodificação da região de memória (RAM ou E/S), que depende dos bits
    mais altos do endereço (2,6~ns);
  \item seleção e extensão do byte ou da meia palavra lida (2,0~ns);
  \item \textit{write-back} e escrita no banco de registradores (1,2~ns).
\end{itemize}

Nenhuma dessas etapas é profunda isoladamente. O atraso vem de executar
todas em série num único ciclo, o que é inerente à organização monociclo.
Ajustes locais, como uma decodificação de região mais rasa ou opções de
otimização da ferramenta, ainda podem render algumas frações de
nanossegundo. Atingir os 100~MHz, porém, exige dividir esse caminho em
estágios, como numa organização em \textit{pipeline}.
